\section{Objetivos de la tesis}
\label{sec:objetivos}

% ~300 palabras. Objetivo general + subobjetivos específicos. Debe quedar
% claro qué se aporta a nivel teórico, metodológico y práctico.

\subsection{Objetivo general}
Evaluar la confiabilidad y la estructura arquitectónica del código producido por un modelo de orquestación multiagente con roles especializados y heterogeneidad completa de modelos frente a un agente único de propósito general en tareas de desarrollo de software asistido por IA.

\subsection{Objetivos específicos}
\begin{enumerate}[label=O\arabic*.]
  \item Revisar el estado del arte sobre orquestación multiagente, especialización de roles y heterogeneidad de modelos en el desarrollo de software asistido por IA, así como las métricas de confiabilidad y estructura arquitectónica empleadas en la literatura.
  \item Definir las métricas de confiabilidad (tasa de errores de compilación y lógicos) y de estructura arquitectónica (modularidad y complejidad cognitiva) que permitan comparar un modelo de orquestación multiagente con roles especializados y heterogeneidad completa de modelos frente a un agente único de propósito general.
  \item Plantear y ejecutar un caso de estudio que permita comparar ambos enfoques en tareas de desarrollo de software asistido por IA.
  \item Evaluar y comparar estadísticamente la confiabilidad y la estructura arquitectónica del código producido por ambos enfoques en el caso de estudio.
\end{enumerate}

\subsection{Aportaciones esperadas}

\begin{description}
  \item[Teóricas y de investigación de base] Resultados previstos de
    claro contenido teórico o que marquen las bases de desarrollos
    posteriores.
  \item[Metodológicas] Resultados de carácter metodológico o que propongan estrategias y técnicas nuevas para el área.
  \item[Prácticas] Desarrollos prácticos previstos que acompañarán a los
    resultados principales de la tesis.
\end{description}
